% This template is dedicated to the public domain under CC0 1.0: use it for
% anything, no credit needed. https://creativecommons.org/publicdomain/zero/1.0/

\documentclass[journal]{IEEEtran}
\usepackage{amsmath}
\usepackage{graphicx}
\usepackage{cite}
\usepackage{maleficium-links}   % violet links; delete for plain IEEEtran
\usepackage{maleficium-footer}   % "Made with Maleficium" mark; delete this line to remove it

\begin{document}

\title{Title of the Paper}
\author{First~Author,~\IEEEmembership{Member,~IEEE,}
        and~Second~Author%
\thanks{Manuscript received on a date to be filled in.}}

\markboth{Journal Name,~Vol.~1, No.~1, Month~Year}%
{Author \MakeLowercase{\textit{et al.}}: Short Title}

\maketitle

\begin{abstract}
The abstract states the problem, the method, the main result and why it
matters, in one paragraph.
\end{abstract}

\begin{IEEEkeywords}
First keyword, second keyword, third keyword.
\end{IEEEkeywords}

\section{Introduction}
\IEEEPARstart{T}{his} template follows IEEE journal conventions. Cite work
with numbered references~\cite{knuth1984, lamport1994}.

\section{Method}
\begin{equation}
  y = W x + b
  \label{eq:model}
\end{equation}
Equation~\eqref{eq:model} is the model.

\section{Results}
Report what you measured.

\section{Conclusion}
Summarise the contribution.

\bibliographystyle{IEEEtran}
\bibliography{references}

\end{document}
